\exercisesfor{3}

\begin{enumerate}
\def\labelenumi{\arabic{enumi}.}
\item
  设 g 为李代数，\(\rho\) 是 g 在 K-模 M 上的表示，\(\sigma\) 是 g 在 M 的张量代数上的相应表示。证明，对称张量子模和斜对称张量子模关于 \(\sigma\) 稳定。
\item
  设 g 为李代数，\(\rho\) 是 g 在 K-模 M 上的表示，\(\sigma\) 是 g 在 Q = L(M, M; M) 上的相应表示。对于 \(f \in Q\) 关于 \(\sigma\)  的不变性充要条件是 \(\rho(x)\) 是关于 f 所定义乘法的 M 的导子。
\item
  设 V 是完美域 K 上的有限维向量空间。
\end{enumerate}

\begin{enumerate}
\def\labelenumi{(\alph{enumi})}
\tightlist
\item
   的恒等表示 \(\mathfrak{gl}(\mathrm{V})\) 在 \(\mathfrak{gl}(\mathrm{V})\) 上规范地定义表示，作用于 \(\bigotimes_{q}^{p}\mathrm{V}\) 和 \(\bigotimes_{q}^{q}\mathrm{V}^{*}\)，从而表示 \(u\mapsto u_q^p\) \(\mathfrak{gl}(\mathrm{V})\) 作用于
\end{enumerate}

\(V_{q}^{p} = \left( \bigotimes_{i=1}^{p} V \right) \otimes \left( \bigotimes_{i=1}^{q} V^{*} \right)\)。证明 \(u_{q}^{p}\) 若 \(u\) 是半单的。（延拓基域，将其归结为 \(K\) 为代数闭域的情形。）证明 \(u_{q}^{p}\) 若 \(u\) 是幂零的。推出，若 \(s\) 与 \(n\) 是 \(u, u_{q}^{p}\) 和 \(n_{q}^{p}\) 的半单与幂零分量，则 \(u_{q}^{p}\) 的半单与幂零分量。

\begin{enumerate}
\def\labelenumi{(\alph{enumi})}
\setcounter{enumi}{1}
\item
  由 (a) 推出，若 \(V_q^p\) 中的一个元素被 \(u_q^p\)  消去，则它也被 \(s_q^p\) 和 \(n_q^p\) 消去。
\item
  由 (b) 和练习 2 推出，若 V 具有不必结合的代数结构且 u 是 V 的导子，则 s 和 n 都是 V 的导子。
\end{enumerate}

\begin{enumerate}
\def\labelenumi{\arabic{enumi}.}
\setcounter{enumi}{3}
\tightlist
\item
  设 K 是一个域。令 g 为李 K-代数。
\end{enumerate}

\begin{enumerate}
\def\labelenumi{(\alph{enumi})}
\item
  设 M、N 为 g-模，其中 N 是 K 上的有限维模。令 \((f_{j})_{1 \leqslant j \leqslant n}\) 为 N 在 K 上的一组基。若存在 M 中的元素 \(e_{1}, \ldots, e_{n}\)，不全为零，使得 \(\sum_{i=1}^{n} e_i \otimes f_i\) 在 \(\mathbf{M} \otimes \mathbf{N}\)  中不变，则由 \(\mathbf{M}_1\) 生成的 \(\mathbf{M}\) 是 \(e_i\)，关于 \(x_{\mathbf{M}}\) 稳定；若进一步 \(\mathbf{N}\) 是单的，则 g 在 \(\mathfrak{g}\) 上的表示 \(\mathbf{M}_1\) 对偶于 g 在 \(\mathfrak{g}\) 上的表示。\(\mathbf{N}\) .
\item
  设 \(M_{1}\) 、\(M_{2}\)、N、\(N^{*}\) 是 K 上的有限维 g-模。设 \(M_{1}\) 与 \(M_{2}\) 是单的，且 g 在 N 和 \(N^{*}\) 上的表示互为对偶。对于 g-模 \(M_{1}\) 要同构于 \(M_{2} \otimes N\)  的一个 g-子模的充要条件是 \(M_{2}\) 同构于 \(M_{1} \otimes N^{*}\)。（使用命题 4；考察 g 在 \(M_{1}^{*} \otimes M_{2} \otimes N\) 上的表示，并将 (a) 应用于后者的对偶表示。）
\end{enumerate}

\begin{enumerate}
\def\labelenumi{\arabic{enumi}.}
\setcounter{enumi}{4}
\tightlist
\item
  设 K 是一个特征为 0 的域。令 V 为有限维 K-向量空间，\(\mathfrak{g} = \mathfrak{sl}(V)\)，U 为 \(\mathfrak{g}\) 的包络代数，\(U_n\) 是 U 中滤子不超过 \(\leqslant n\) 的元素集合；\(U^n\) 是 U 中阶数为 \(n\) 的、在 \(\mathfrak{g}\) 上定义的齐次对称张量的像，且\((\alpha_1, \ldots, \alpha_m)\) 是向量空间 g 的一组基。令 \(W_{s} = \bigotimes^{s} V\)。g 的恒等表示在 \(W_{s}\) 上定义了 g 的一个表示，并将其延拓为同态 \(\pi_s\)，从 U 到代数 \(M_s = \mathcal{L}(W_s)\)。
\end{enumerate}

\begin{enumerate}
\def\labelenumi{(\alph{enumi})}
\tightlist
\item
  令\(M_{s}^{\prime}\)为由\(M_{s}\)中形如\(\beta_{1} \otimes \beta_{2} \otimes \cdots \otimes \beta_{s}\)的元素生成的子空间，其中\(\beta_{i} \in \mathcal{L}(V)\)，且\(\beta_{i} = 1\)对至少一个 i 等于 1。令 z 为\(U_{s}\)中的一个元素，并记其\(z_{0} = \sum_{1 \leqslant i_{1}, \ldots, i_{s} \leqslant m} a_{i_{1} \ldots i_{s}} \alpha_{i_{1}} \ldots \alpha_{i_{s}}\)为\(U^{s}\)中的分量；其中\(a_{i_{1} \ldots i_{s}}\)关于指标置换是对称的。证明
\end{enumerate}

\[
\pi_ {s} (z) \equiv s! \sum_ {1 \leqslant i _ {1}, \dots , i _ {s} \leqslant m} a _ {i _ {1} \dots i _ {s}} \alpha_ {i _ {1}} \otimes \dots \otimes \alpha_ {i _ {s}} \pmod {\mathrm{M} _ {s} ^ {\prime}}.
\]

\begin{enumerate}
\def\labelenumi{(\alph{enumi})}
\setcounter{enumi}{1}
\tightlist
\item
  由 (a) 推出，对任意\(z \in \mathbf{U}\)，存在一个有限维表示\(\pi\)，其作用于\(\mathbf{U}\)满足\(\pi(z) \neq 0\)不为 0。（关于本题的后续内容，参见 § 7，习题 3。）
\end{enumerate}

\begin{enumerate}
\def\labelenumi{\arabic{enumi}.}
\setcounter{enumi}{5}
\tightlist
\item
 设\(\mathbf{K}\)是一个域。令\(\mathfrak{g}\)为李 K-代数，\(x \mapsto x_{\mathbf{M}}\)是\(\mathfrak{g}\)的有限维表示，\((e_1, \ldots, e_n)\)是\(\mathbf{M}\)的一组基，而\((e_1^*, \ldots, e_n^*)\)是对偶基。
\end{enumerate}

\begin{enumerate}
\def\labelenumi{(\alph{enumi})}
\tightlist
\item
  若\(f\)是定义在\(\bigotimes^{r} \mathbf{M}\)上的线性形式，则：
\end{enumerate}

\[
f = \sum_ {1 \leqslant i _ {1} \dots , i _ {r} \leqslant n} f (e _ {i _ {1}} \otimes \dots \otimes e _ {i _ {r}}) (e _ {i _ {1}} ^ {*} \otimes \dots \otimes e _ {i _ {r}} ^ {*}).
\]

\begin{enumerate}
\def\labelenumi{(\alph{enumi})}
\setcounter{enumi}{1}
\tightlist
\item
  令\(\beta\)为定义在\(\mathbf{M}\)上且在\(\mathfrak{g}\)下不变的非退化双线性形式。令\((e_1', \ldots, e_n')\)为\(\mathbf{M}\)的一组基，满足\(\beta(e_i, e_j') = \delta_{ij}\)。令\(f\)为定义在\(\bigotimes^r \mathbf{M}\)上的不变线性形式。由 (a) 推出：
\end{enumerate}

\[
\sum_ {1 \leqslant i _ {1}, \dots , i _ {r} \leqslant n} f (e _ {i _ {1}} \otimes \dots \otimes e _ {i _ {r}}) (e _ {i _ {1}} ^ {\prime} \otimes \dots \otimes e _ {i _ {r}} ^ {\prime})
\]

是\(\otimes\) M 中的一个不变元，且与基\((e_i)\)的选择无关。

\begin{enumerate}
\def\labelenumi{(\alph{enumi})}
\setcounter{enumi}{2}
\item
  令 U 为 g 的包络代数，U* 为 U 的对偶空间。g 的伴随表示可延拓为表示\(x \mapsto x_{\mathrm{U}}\)：它在 U 上由\(x_{\mathrm{U}}u = xu - ux\)定义，对所有\(u \in U\)均成立。因此 U* 具有 g-模结构。对于 U* 中的元素 f，其为不变元的充要条件是\(f(uv) = f(vu)\)，对 U 中所有 u、v 均成立。
\item
  令\(\mathfrak{h}\)为\(\mathfrak{g}\)的有限维理想，\(\gamma\)为定义在\(\mathfrak{g}\)上的不变双线性形式，且其在\(\mathfrak{h}\)上的限制非退化。令\((y_i)_{1\leq i\leq n}, (y_j')_{1\leq j\leq n}\)为\(\mathfrak{h}\)的两组基，满足\(\gamma(y_i,y_j') = \delta_{ij}\)。令\(r\)为一个\(\geqslant 1\)的整数。令\(f\)为 U 上的线性形式，满足\(f(w) = f(vu)\)，对 U 中所有\(u,v\)均成立。由 (b) 和 (c) 推出：
\end{enumerate}

\[
\sum_ {1 \leqslant t _ {1}, \dots t _ {r} \leqslant n} f (y _ {t _ {1}} y _ {t _ {2}} \cdot \cdot \cdot y _ {t _ {r}}) y _ {t _ {1}} ^ {\prime} y _ {t _ {2}} ^ {\prime} \cdot \cdot \cdot y _ {t _ {r}} ^ {\prime}
\]

是 U 中心里与基\((y_{i})\)的选择无关的元素。特别地，这给出了 Casimir 元。

\begin{enumerate}
\def\labelenumi{(\alph{enumi})}
\setcounter{enumi}{4}
\tightlist
\item
  令\(\theta\)为\(\{1, \ldots, r\}\)的一个置换。用类似的论证证明：
\end{enumerate}

\[
\sum_ {1 \leqslant i _ {1}, \dots i _ {r} \leqslant n} f (y _ {i _ {\theta (1)}} y _ {i _ {\theta (2)}} \dots y _ {i _ {\theta (r)}}) y _ {i _ {1}} ^ {\prime} y _ {i _ {2}} ^ {\prime} \dots y _ {i _ {r}} ^ {\prime}
\]

是 U 中心里与基\((y_{i})\)的选择无关的元素。

\begin{enumerate}
\def\labelenumi{\arabic{enumi}.}
\setcounter{enumi}{6}
\item
  令 g 为复李代数，\(\beta\)为其 Killing 形式，\(g_{0}\)为将基域限制为 R 后由 g 得到的李代数。证明\(g_{0}\)的 Killing 形式是\(\beta\)的实部的两倍。
\item
  令\(\mathfrak{g}\)为实李代数。多线性形式
\end{enumerate}

\[
(x _ {1}, \dots , x _ {n}) \mapsto \operatorname{Tr} (\operatorname{ad} x _ {1} \circ \operatorname{ad} x _ {2} \circ \dots \circ \operatorname{ad} x _ {n})
\]

在\(\mathfrak{g}^n\)上是不变的。

\begin{enumerate}
\def\labelenumi{\arabic{enumi}.}
\setcounter{enumi}{8}
\item
  令 g 为李 K-代数，\(\rho\)和\(\rho'\)是 g 在 K-模 V、\(V'\)上的半单表示，且\(\phi\)是从 g-模 V 到 g-模\(V'\)的满同态。证明，\(\phi\)将 V 的不变元集合映到\(V'\)的不变元集合（使用命题 6）。
\item
  设 K 是一个域。令 g 为李 K-代数，\(M_{1}\)、\(M_{2}\)是两个不同构的有限维简单 g-模。若\(K_{1}\)是 K 的可分扩张，证明不存在简单的\(g_{(K_{1})}\)-模，它同时同构于\(g_{(K_{1})}\)-模\(M_{1(K_{1})}\)的一个子模，也同构于\(g_{(K_{1})}\)-模\(M_{2(K_{1})}\)的一个子模。（使用命题 4，并注意：\(\neq0\)在\(M_{1(K_{1})}^{*}\otimes_{K}M_{2(K_{1})}\)中表示一个不变元，意味着\(\neq0\)在\(M_{1}^{*}\otimes_{K}M_{2}\)中表示一个不变元（《代数》第二章 § 5，第 3 号，定理 1）。）
\item
  令\(\mathfrak{w}\)为 § 1 习题 21 中定义的李 \(p\)-代数。证明\(\mathfrak{w}\)的 Killing 形式为零。
\end{enumerate}

¶ 12. 设 K 是一个域。令 g 为李 K-代数，M 为 g-模。令\(C^{p}(\mathfrak{g}, \mathrm{M})\)表示从\(g^{p}\)到 M 的交错线性映射空间。记\(C^{0}(\mathfrak{g}, \mathrm{M}) = \mathrm{M}\)，并且当 p\textless{} 0 时，\(C^{p}(\mathfrak{g}, \mathrm{M}) = \{0\}\)。令\(C^{*}(\mathfrak{g}, \mathrm{M})\)为各个\(C^{p}(\mathfrak{g}, \mathrm{M})\)的直和。\(C^{*}(\mathfrak{g}, \mathrm{M})\)中的元素称为 g 取值于 M 的上链；\(C^{p}(\mathfrak{g}, \mathrm{M})\)中的元素称为 p 次上链。对于所有\(y \in g\)，记\(i(y)\)为\(C^{*}(\mathfrak{g}, \mathrm{M})\)的自同态；它将每个子空间\(C^{p}(\mathfrak{g}, \mathrm{M})\)映到\(C^{p-1}(\mathfrak{g}, \mathrm{M})\)，并且当 p\textgreater{} 0 时由下式给出：

\[
(i (y) f) \left(x _ {1}, \dots , x _ {p - 1}\right) = f \left(y, x _ {1}, \dots , x _ {p - 1}\right).\tag{1}
\]

有\(i(y)^{2}=0\)。

\begin{enumerate}
\def\labelenumi{(\alph{enumi})}
\tightlist
\item
  g 的伴随表示与 g 在 M 上的表示在从\(g^{p}\)到 M 的多线性映射空间上定义了 g 的一个表示。证明\(C^{p}(g, M)\)在此表示下稳定。令\(\theta\)为 g 在\(C^{*}(g, M)\)上的上述表示。证明：
\end{enumerate}

\[
\theta (x) i (y) - i (y) \theta (x) = i ([ x, y ])
\]

对 g 中所有 x、y。

\begin{enumerate}
\def\labelenumi{(\alph{enumi})}
\setcounter{enumi}{1}
\tightlist
\item
  证明存在唯一的自同态\(d\)，它作用于\(\mathrm{C}^*(\mathfrak{g},\mathrm{M})\)，将\(\mathrm{C}^p (\mathfrak{g},\mathrm{M})\)映到\(\mathrm{C}^{p + 1}(\mathfrak{g},\mathrm{M})\)，使得
\end{enumerate}

\[
d i (y) + i (y) d = \theta (y)\tag{3}
\]

对所有\(y \in \mathfrak{g}\)。（按上链次数作归纳。）证明，对\(f \in \mathrm{C}^p(\mathfrak{g}, \mathbf{M})\)：

\[
\begin{array}{r l} d f (x _ {1}, x _ {2}, \dots , x _ {p + 1}) & = \sum_ {i <   j} (- 1) ^ {i + j} f ([ x _ {i}, x _ {j} ], x _ {1}, \dots , \hat {x} _ {i}, \dots , \hat {x} _ {j}, \dots , x _ {p + 1}) \\ & \quad + \sum_ {i} (- 1) ^ {i + 1} (x _ {i}) _ {\mathrm{M}} f (x _ {1}, \dots , \hat {x} _ {i}, \dots , x _ {p + 1}) \end{array}
\]

（其中字母上方的符号\^{}表示省略该字母）。

\begin{enumerate}
\def\labelenumi{(\alph{enumi})}
\setcounter{enumi}{2}
\tightlist
\item
  证明，对所有\(y \in g\)，
\end{enumerate}

\[
d \theta (y) = \theta (y) d.\tag{4}
\]

（先利用 (2) 和 (3) 证明\(d\theta(y) - \theta(y)d\)与每个\(i(x)\)交换。然后按上链次数作归纳。）

\begin{enumerate}
\def\labelenumi{(\alph{enumi})}
\setcounter{enumi}{3}
\item
  证明\(d^2 = 0\)。（先利用 (3) 和 (4) 证明\(d^2\)与每个\(i(x)\)交换。然后按上链次数作归纳。）
\item
  \(d\)在\(\mathrm{C}^p (\mathfrak{g},\mathrm{M})\)上的限制的核为\(\mathrm{Z}^p (\mathfrak{g},\mathrm{M})\)，其中元素称为取值于 M 的\(p\)次上循环。\(d\)在\(\mathrm{C}^{p - 1}(\mathfrak{g},\mathrm{M})\)上的限制的像为\(\mathrm{B}^p (\mathfrak{g},\mathrm{M})\)，其中元素称为取值于 M 的\(p\)次上余边界。因此\(\mathrm{B}^p (\mathfrak{g},\mathrm{M})\subset \mathrm{Z}^p (\mathfrak{g},\mathrm{M})\)。商空间
\end{enumerate}

\[
\mathrm{Z} ^ {p} (\mathfrak {g}, \mathrm{M}) / \mathrm{B} ^ {p} (\mathfrak {g}, \mathrm{M}) = \mathrm{H} ^ {p} (\mathfrak {g}, \mathrm{M})
\]

称为 g 的 p 次、取值于 M 的上同调空间。\(\mathrm{H}^{p}(\mathfrak{g},\mathrm{M})\)的直和记为\(\mathrm{H}^{*}(\mathfrak{g},\mathrm{M})\)。证明\(\mathrm{H}^{0}(\mathfrak{g},\mathrm{M})\)可与 M 的不变元集合等同。令\(\phi\)为从 g-模 M 到 g-模 N 的同态。对所有\(f \in \mathrm{C}^{p}(\mathfrak{g},\mathrm{M})\)，\(\phi \circ f \in \mathrm{C}^{p}(\mathfrak{g},\mathrm{N})\)，因此\(\phi\)可延拓为 K-线性映射\(\phi'\)：\(\mathrm{C}^{*}(\mathfrak{g},\mathrm{M}) \to \mathrm{C}^{*}(\mathfrak{g},\mathrm{N})\)。证明\(\phi' \circ d = d \circ \phi'\)。由此推出\(\phi'\)定义了同态

\[
\tilde {\phi} \colon \mathrm{H} ^ {*} (\mathfrak {g}, \mathrm{M}) \to \mathrm{H} ^ {*} (\mathfrak {g}, \mathrm{N})
\]

其称为与\(\phi\)相关的同态。

\begin{enumerate}
\def\labelenumi{(\alph{enumi})}
\setcounter{enumi}{5}
\tightlist
\item
  令 L 为 M 的子 g-模，N 为商 g-模 M/L。规范同态\(L \stackrel{\iota}{\rightarrow} M \stackrel{p}{\rightarrow} N\)定义了同态：
\end{enumerate}

\[
\begin{array}{l} \mathrm{C} ^ {*} (\mathfrak {g}, \mathrm{L}) \xrightarrow {\iota^ {\prime}} \mathrm{C} ^ {*} (\mathfrak {g}, \mathrm{M}) \xrightarrow {p ^ {\prime}} \mathrm{C} ^ {*} (\mathfrak {g}, \mathrm{N}), \\ \mathrm{H} ^ {n} (\mathfrak {g}, \mathrm{L}) \xrightarrow {\iota^ {n}} \mathrm{H} ^ {n} (\mathfrak {g}, \mathrm{M}) \xrightarrow {p ^ {n}} \mathrm{H} ^ {n} (\mathfrak {g}, \mathrm{N}). \end{array}
\]

证明\(\iota'\)是单射，且其像为\(p'\)的核，并且\(p'\)是满射。对任意\(c \in \mathrm{H}^n(\mathfrak{g},\mathrm{N})\)，令\(z \in \mathrm{Z}^n(\mathfrak{g},\mathrm{N})\)为\(c\)的一个代表元，并令\(a \in \mathrm{C}^n(\mathfrak{g},\mathrm{M})\)满足\(p'(a) = z\)；证明\(da \in \mathrm{Z}^{n+1}(\mathfrak{g},\mathrm{L})\)，且其在\(\mathrm{H}^{n+1}(\mathfrak{g},\mathrm{L})\)中的类只依赖于\(c\)；将此类记为\(\delta^n c\)。证明以下同态序列：

\[
\begin{array}{r l} 0 \longrightarrow & \mathrm{H} ^ {0} (\mathfrak {g}, \mathrm{L}) \stackrel {{\iota ^ {0}}} {{\longrightarrow}} \mathrm{H} ^ {0} (\mathfrak {g}, \mathrm{M}) \stackrel {{p ^ {0}}} {{\longrightarrow}} \mathrm{H} ^ {0} (\mathfrak {g}, \mathrm{N}) \stackrel {{\delta^ {0}}} {{\longrightarrow}} \\ & \longrightarrow \mathrm{H} ^ {1} (\mathfrak {g}, \mathrm{L}) \stackrel {{\iota ^ {1}}} {{\longrightarrow}} \mathrm{H} ^ {1} (\mathfrak {g}, \mathrm{M}) \stackrel {{p ^ {1}}} {{\longrightarrow}} \mathrm{H} ^ {1} (\mathfrak {g}, \mathrm{N}) \stackrel {{\delta^ {1}}} {{\longrightarrow}} \dots \end{array}
\]

是一个正合序列。

\begin{enumerate}
\def\labelenumi{(\alph{enumi})}
\setcounter{enumi}{6}
\tightlist
\item
  按 (f) 的记号，正合序列：
\end{enumerate}

\[
0 \to \mathscr {L} (\mathrm{N}, \mathrm{L}) \to \mathscr {L} (\mathrm{N}, \mathrm{M}) \to \mathscr {L} (\mathrm{N}, \mathrm{N}) \to 0
\]

定义了一个正合序列

\[
\mathrm{H} ^ {0} (\mathfrak {g}, \mathscr {L} (\mathrm{N}, \mathrm{M})) \rightarrow \mathrm{H} ^ {0} (\mathfrak {g}, \mathscr {L} (\mathrm{N}, \mathrm{N})) \rightarrow \mathrm{H} ^ {1} (\mathfrak {g}, \mathscr {L} (\mathrm{N}, \mathrm{L})).
\]

 N 的恒等映射是\(\mathcal{L}(\mathrm{N},\mathrm{N})\)中的一个不变元，因而是\(\mathrm{H}^{0}(\mathfrak{g},\mathcal{L}(\mathrm{N},\mathrm{N}))\)中的一个元素；令 c 为它在\(\mathrm{H}^{1}(\mathfrak{g},\mathcal{L}(\mathrm{N},\mathrm{L}))\)中的像。于是，M 中存在补充 L 的子 g-模的充要条件是 c=0。（此条件意味着 u 是\(\mathrm{H}^{0}(\mathfrak{g},\mathcal{L}(\mathrm{N},\mathrm{M}))\)中某个元素的像，即从 g-模 N 到 g-模 M 的一个同态；此时\(v(\mathrm{N})\)就是所求的补子。）

\begin{enumerate}
\def\labelenumi{(\alph{enumi})}
\setcounter{enumi}{7}
\item
  证明\(Z^{1}(\mathfrak{g},\mathfrak{g})\)（其中 g 通过伴随表示视为 g-模）可与 g 的导子空间等同，且\(B^{1}(\mathfrak{g},\mathfrak{g})\)可与 g 的内导子空间等同。
\item
  令\(\mathfrak{a}\)和\(\mathfrak{b}\)为两个李 K-代数，其中\(\mathfrak{b}\)为交换代数，且\(\mathfrak{b} \xrightarrow{\lambda} \mathfrak{g} \xrightarrow{\mu} \mathfrak{a}\)是\(\mathfrak{a}\)通过\(\mathfrak{b}\)的扩张。对于所有\(x \in \mathfrak{g}\)，\(\mathrm{ad}_{\mathfrak{a}} x\)在\(\mathfrak{b}\)上的限制只取决于\(x\)模\(\mathfrak{b}\)的类，因此\(\mathfrak{a}\)上具有\(\mathfrak{b}\)-模结构。令\(\nu\)为从\(\mathfrak{a}\)到\(\mathfrak{g}\)的 K-线性映射，满足\(\mu \circ \nu\)是\(\mathfrak{a}\)的恒等映射。对于\(x, y\)中的\(\mathfrak{a}\)，记\(f(x, y) = [\nu x, \nu y] - \nu([x, y])\)。证明\(f \in Z^2(\mathfrak{a}, \mathfrak{b})\)，且\(c\)（即\(f\)在\(H^2(\mathfrak{a}, \mathfrak{b})\)中的类）不依赖于\(\nu\)的选择。对于\(\mathfrak{a}\)通过\(\mathfrak{b}\)的两个扩张，若它们定义了相同的\(\mathfrak{a}\)-模结构（作用于\(\mathfrak{b}\)），则二者等价的充要条件是相应的类\(c \in H^2(\mathfrak{a}, \mathfrak{b})\)相同。该扩张无本质的充要条件是\(c = 0\)。若 B 是\(\mathfrak{a}\)-模且\(c \in H^2(\mathfrak{a}, \mathfrak{B})\)，则存在\(\mathfrak{a}\)通过 B（将 B 视为交换李代数）的扩张，它在 B 上定义给定的\(\mathfrak{a}\)-模结构，并实现\(c\)这一\(H^2(\mathfrak{a}, \mathfrak{B})\)中的给定元素。
\item
  令 g 为定义在 K 上的有限维李代数，U 为其包络代数，h 为 g 的理想，β 为定义在 g 上且在 h 上限制非退化的不变双线性形式，\((y_{i})_{1 \leqslant i \leqslant n}\)和\((y_{i}')_{1 \leqslant i \leqslant n}\)为 h 的两组基，满足
\end{enumerate}

 \(\beta(y_i, y_i') = \delta_{ij}, t = \sum_{i=1}^{n} y_i y_i' \in \mathbf{U}, \mathbf{M}\)，\(g\)-模中的\(\rho\)是\(\mathrm{C}^*(\mathfrak{g}, \mathbf{M})\)的自同态，它将每个\(\mathrm{C}^p(\mathfrak{g}, \mathbf{M})\)映到\(\mathrm{C}^{p-1}(\mathfrak{g}, \mathbf{M})\)，并且当\(p > 0\)时由下式定义：

\[
(\wp f) (x _ {1}, \dots , x _ {p - 1}) = \sum_ {k = 1} ^ {n} (y _ {k}) _ {\mathrm{M}} f (y _ {k} ^ {\prime}, x _ {1}, \dots , x _ {p - 1}).
\]

最后，令\(\Gamma\)为\(C^*(g, M)\)的自同态，它延拓\(t_M\)，将每个上链\(f\)（次数为\(>0\)）映到上链\(t_M \circ f\)。证明\(\rho d + d\rho = \Gamma\)。（若对\(x \in g\)，写作：

\[
[ x, y _ {k} ] = \sum_ {l = 1} ^ {n} a _ {k l} y _ {l}, \quad [ x, y _ {k} ^ {\prime} ] = \sum_ {l = 1} ^ {n} a _ {k l} ^ {\prime} y _ {l} ^ {\prime},
\]

·先证明\(a_{kl} = -a_{lk}^{\prime}\)。)由此推出\(\Gamma d = d\Gamma\)。若 M 是 K 上有限维的简单模，\(\beta\)是相应的双线性形式，且 dim\(\mathfrak{h}\)不被 K 的特征整除，证明对所有 p 都有\(\mathrm{H}^{p}(\mathfrak{g},\mathrm{M})=\{0\}\)。~（利用命题 12，证明\(\Gamma\)是\(\mathrm{C}^{*}(\mathfrak{g},\mathrm{M})\)的自同构，因而在\(\mathrm{Z}^{p}(\mathfrak{g},\mathrm{M})\)和\(\mathrm{B}^{p}(\mathfrak{g},\mathrm{M})\)上诱导自同构。对于

\[
f \in \mathrm{Z} ^ {p} (\mathfrak {g}, \mathrm{M}), \quad \Gamma f = (d \rho + \rho d) f = d \rho f \in \mathrm{B} ^ {p} (\mathfrak {g}, \mathrm{M}),
\]

于是\(Z^{p}(\mathfrak{g},\mathbf{M}) = \mathrm{B}^{p}(\mathfrak{g},\mathbf{M}).\)

